% =============================================================================
% Abstract and Introduction
% Author: Abhishek Kumar Patel (2401010020)
% Included from docs/report/main.tex
% =============================================================================

\begin{abstract}
Pipelining is the basic technique that allows a processor to work on several instructions at the
same time and thereby raise instruction throughput without raising the clock frequency. This
project presents the design and implementation of a 32-bit, five-stage pipelined processor for a
subset of the open RISC-V RV32I base integer instruction set, written in Verilog HDL. The datapath
is divided into Instruction Fetch (IF), Instruction Decode (ID), Execute (EX), Memory Access (MEM)
and Write-Back (WB) stages separated by four pipeline registers. The supported subset covers
register--register and register--immediate arithmetic, logic, shift and comparison instructions,
word loads and stores, conditional branches, jumps and \texttt{LUI}. Data hazards are resolved by a
forwarding unit that routes results from the EX/MEM and MEM/WB registers to the ALU inputs, by a
register file with write-first bypass, and by a hazard detection unit that inserts a one-cycle bubble
for load-use dependencies. Control hazards are handled by resolving branches and jumps in the EX stage
with a predict-not-taken policy and flushing the two wrongly fetched instructions. Every module is
verified with a self-checking testbench in Icarus Verilog, and complete programs are verified at
system level with waveform inspection in GTKWave. This report describes the Phase-I design and
implementation progress.
\end{abstract}

\begin{IEEEkeywords}
RISC-V, RV32I, pipelined processor, Verilog HDL, data forwarding, hazard detection, computer architecture
\end{IEEEkeywords}

\section{Introduction}

RISC-V is an open, royalty-free instruction set architecture (ISA) originally developed at the
University of California, Berkeley~\cite{riscv_spec_unpriv}. Its small and regular base integer
instruction set, fixed 32-bit instruction length and a small number of instruction formats make it
well suited for teaching and for implementing processors in hardware description languages.

A single-cycle processor completes every instruction in one long clock cycle, so the clock period
is set by the slowest instruction and most of the hardware is idle for most of the cycle.
Pipelining divides instruction execution into stages and starts a new instruction in every clock
cycle, so that, in the ideal case, one instruction completes per cycle and throughput increases by
up to the number of stages~\cite{patterson_riscv, hennessy_quant}. The cost of pipelining is the
appearance of \emph{hazards}: an instruction may need a result that an earlier instruction has not
yet written back (data hazard), or the processor may fetch instructions from the wrong path before a
branch is resolved (control hazard). Correct and efficient handling of these hazards is the central
design problem of a pipelined processor.

The objectives of this project are:
\begin{enumerate}
  \item to design a five-stage (IF, ID, EX, MEM, WB) pipelined datapath and control unit for a
        subset of the RV32I instruction set;
  \item to resolve data hazards with operand forwarding, a register-file bypass and a load-use
        stall, and control hazards with branch resolution in EX and pipeline flushing;
  \item to implement the design as modular, synthesizable-style Verilog with a clearly defined
        interface between modules so that it can be developed in parallel by a team; and
  \item to verify every module and the complete processor through self-checking simulation.
\end{enumerate}

The design was studied alongside an existing open-source pipelined RISC-V implementation~\cite{protocolx_repo};
all modules in this project were written by the team against a common interface specification.

The rest of this report is organised as follows. Section~II reviews related literature.
Section~III presents the proposed solution: the overall architecture, each pipeline stage, hazard
handling and the verification methodology. Section~IV lists the individual contributions of the team
members, followed by the references.
